\subsection{Experiments with real-world functions}

We finally optimize different real-world cost functions. The 
first one (NNet) returns the predictive accuracy of a neural network on a random train/test partition of 
the Boston Housing dataset \cite{BacheLichman:2013}. The variables to optimize are the weight-decay parameter and the number of training
iterations for the neural network. The second function (Hydrogen) returns the amount of
hydrogen production of a particular bacteria in terms of the PH and Nitrogen levels of the growth medium \cite{Burrows:2009}.
The third one (Portfolio) returns the ratio of the mean and the standard deviation (the Sharpe ratio) of the 
1-year ahead returns generated by simulations from a multivariate time-series model that is adjusted to the daily returns of
stocks AXP, BA and HD. The time-series model is formed by univariate GARCH models connected with a Student's $t$ copula \cite{jondeau:2006}.
These three functions (NNet, Hydrogen and Portfolio) have as domain $[0,1]^2$. 
Furthermore, in these examples, the ground truth function that we want to optimize is unknown and is only available 
through noisy measurements. To obtain a ground truth, we approximate each cost function as the
predictive distribution of a GP that is adjusted to data sampled from the original function (1000 uniform samples for NNet and Portfolio
and all the available data for Hydrogen \cite{Burrows:2009}).
Finally, we also consider another real-world function that returns the walking
speed of a bipedal robot \cite{Westervelt:2007}.
This function is defined in $[0,1]^8$ and its inputs are the parameters
of the robot's controller. In this case the ground truth function is noiseless
and can be exactly evaluated through expensive numerical simulation. We consider two versions of this problem (Walker A) with zero-mean, additive noise of $\sigma=0.01$ and (Walker B) with $\sigma=0.1$.

Figure \ref{fig:resultsNonAnalyticRealWorldFunctions} shows the median IR values
obtained by each method on each function across 250 random initializations, except in Hydrogen where we used 500 due to its higher level of noise.
Overall, PES, ES and PES-NB perform similarly in NNet, Hydrogen and Portfolio.
EI performs rather poorly in these first three functions. 
This method seems to make excessively greedy decisions and fails to explore the search space enough.
This strategy seems to be advantageous in Walker A, where EI obtains the best results.
By contrast, PES, ES and PES-NB tend to explore more in this latter dataset. This leads to worse results than those of EI.
Nevertheless, PES is significantly better than PES-NB and ES in both Walker datasets and better than EI in the noisier Walker B.
In this case, the fully Bayesian treatment of hyper-parameters performed by PES produces improvements in performance.

\begin{figure}
\begin{center}
\includegraphics[width=1.00\linewidth]{figures/experimentNonAnalyticRealWorldFunctions/figure.pdf}
\end{center}
\vspace{-0.5em}
\caption{\small Median of the immediate regret (IR) for the methods PES, PES-NB, ES and EI in the experiments with non-analytic
real-world cost functions.
}\label{fig:resultsNonAnalyticRealWorldFunctions}
\vspace{-0.5em}
\end{figure}
